\section{An exact-cover reduction}\label{sec:core}

Let $S_1,\ldots,S_n\subseteq[p]$, where $p\geq2$, and let
$A\in\{0,1\}^{p\times n}$ be their incidence matrix.  An exact cover is a
vector $x\in\{0,1\}^n$ satisfying $Ax=\one_p$.  Set
$W=\{1,\ldots,n,g\}$, where $g$ is an additional index.  Choose a rational
number $\eta^2$ such that
\begin{equation}\label{eq:eta-window}
  \max\left\{\frac p{12},\frac{p-2}{4}\right\}
  \leq\eta^2<\frac p4,
\end{equation}
and define the symmetric matrix $M$ on $W$ by
\begin{equation}\label{eq:cover-matrix}
  M_{ij}=4\mathbf 1_{i=j}+\abs{S_i\cap S_j},\qquad
  M_{ig}=\frac{\abs{S_i}}2,\qquad
  M_{gg}=\frac p4+\eta^2.
\end{equation}
The interval in \eqref{eq:eta-window} is nonempty for every $p\geq2$;
the choice $\eta^2=(p-1)/4$ belongs to it.

\begin{theorem}[Exact-cover encoding]\label{thm:exact-cover}
The matrix $M$ in \eqref{eq:cover-matrix} is positive definite.  For
every integer vector $z$, the following equivalence holds:
\begin{equation}\label{eq:cover-violators}
  g_M(z)<0
  \quad\Longleftrightarrow\quad
  z=(x,-1),\quad x\in\{0,1\}^n,\quad Ax=\one_p.
\end{equation}
Every such vector has
\[
  g_M(z)=-\left(\frac p2-2\eta^2\right)\in[-1,0).
\]
\end{theorem}

\begin{proof}
Take $\eta=\sqrt{\eta^2}>0$.  In $\R^{n+p+1}$, define
\[
  u_i=(2e_i,A_i,0),\qquad
  u_g=(\zero_n,\tfrac12\one_p,\eta),
\]
where $A_i$ is column $i$ of $A$.  Their Gram matrix is $M$.  They are
linearly independent: the first block forces every coefficient of the
$u_i$ to vanish, and the last coordinate then forces the coefficient of
$u_g$ to vanish.  Thus $M\succ0$.

For $z=(x,k)\in\Z^n\times\Z$ and $y=Ax$, expansion gives
\begin{equation}\label{eq:cover-slack}
  g_M(x,k)
  =4\sum_i x_i(x_i-1)
   +\sum_e\bigl(y_e^2+(k-1)y_e\bigr)
   +k(k-1)\left(\frac p4+\eta^2\right).
\end{equation}
Indeed, the quadratic part is
\[
  4\norm{x}^2+\norm{y+\tfrac k2\one_p}^2+k^2\eta^2,
\]
and the linear part is
$\sum_i(4+\abs{S_i})x_i+k(p/4+\eta^2)$.  The equality
$\sum_i\abs{S_i}x_i=\sum_e y_e$ follows from the incidence matrix.

Every term $x_i(x_i-1)$ is nonnegative on the integers.  For $k=0$, the
element terms in \eqref{eq:cover-slack} are $y_e(y_e-1)\geq0$; for $k=1$,
they are $y_e^2\geq0$.  For $k\geq2$, their real minima give
\[
  g_M(x,k)
  \geq-\frac{p(k-1)^2}{4}
       +k(k-1)\left(\frac p4+\eta^2\right)
  =(k-1)\left(\frac p4+k\eta^2\right)>0.
\]

Next let $k=-t$ with $t\geq2$.  The minimum of
$y^2-(t+1)y$ over integer $y$ is
$-\lfloor(t+1)^2/4\rfloor$.  If $t$ is odd, then $t\geq3$ and
\[
  g_M(x,-t)
  \geq-\frac{p(t+1)^2}{4}
       +t(t+1)\left(\frac p4+\eta^2\right)
  =(t+1)\left(t\eta^2-\frac p4\right)\geq0.
\]
If $t$ is even, then
\[
  g_M(x,-t)
  \geq-\frac{pt(t+2)}4
       +t(t+1)\left(\frac p4+\eta^2\right)
  =t\left((t+1)\eta^2-\frac p4\right)\geq0.
\]
Both last inequalities follow from $3\eta^2\geq p/4$.

Finally, at $k=-1$, \eqref{eq:cover-slack} becomes
\begin{equation}\label{eq:cover-layer-minus-one}
  g_M(x,-1)
  =4\sum_i x_i(x_i-1)+\norm{Ax-\one_p}^2
    -\left(\frac p2-2\eta^2\right).
\end{equation}
The last parenthesis belongs to $(0,1]$ by
\eqref{eq:eta-window}.  The other two summands are nonnegative integers.
Their sum is strictly smaller than that parenthesis exactly when both
vanish, which is equivalent to $x\in\{0,1\}^n$ and $Ax=\one_p$.
This proves both directions of \eqref{eq:cover-violators} and the claimed
slack.
\end{proof}

The construction uses the incidence identity
$\norm{A_i}^2=\one_p\trans A_i$.  Geometrically, every incidence column
lies on the sphere through the origin centered at $\one_p/2$.
The first block of the generators penalizes integer coordinates outside
$\{0,1\}$, and the distinguished generator selects the exact-cover
equations in the layer $k=-1$.  The realization may use an irrational
$\eta$; only the rational entries of $M$ are part of the output.

For the rest of this section, fix
$\eta^2=(p-1)/4$.  The distance vector $d=d(M)$ is then
\begin{equation}\label{eq:cover-distances}
\begin{aligned}
  d_{0i}&=4+\abs{S_i}, &
  d_{0g}&=\frac{2p-1}{4},\\
  d_{ij}&=8+\abs{S_i\symdiff S_j}\quad(i\ne j), &
  d_{ig}&=\frac{2p+15}{4}.
\end{aligned}
\end{equation}
In particular, $4d$ is integral with positive entries at most $4p+32$.
The identity \eqref{eq:hypermetric-g} and \cref{thm:exact-cover} show that
the complete set of hypermetric violators is
\begin{equation}\label{eq:cover-hypermetric-vectors}
  b_T=(2-\abs{T})e_0+\sum_{i\in T}e_i-e_g,
  \qquad T\text{ an exact cover},
\end{equation}
and each satisfies $Q(b_T,d)=1/2$.

To place the reduction inside the standard relaxations, write
\begin{equation}\label{eq:core-scaling}
  s=\delta\trans M^{-1}\delta,\qquad
  N=8n+2p+1,\qquad \varepsilon=\frac1N,
\end{equation}
and define the binary moment matrix
\begin{equation}\label{eq:core-moment}
  Y_\varepsilon=
  \begin{pmatrix}
    1&\varepsilon\delta\trans\\
    \varepsilon\delta&\varepsilon M
  \end{pmatrix}.
\end{equation}

\begin{lemma}[Polynomial scaling]\label{lem:scaling}
For every exact-cover instance with $p\geq2$,
\begin{equation}\label{eq:scaling-bound}
  \max_{i\in W}M_{ii}
  \leq s\leq4n+p+\frac1{4(p-1)}<\frac N2.
\end{equation}
The matrix $Y_\varepsilon$ satisfies
$(Y_\varepsilon)_{ii}=(Y_\varepsilon)_{0i}$ and
\begin{equation}\label{eq:core-definiteness-margin}
  Y_\varepsilon-\theta e_0e_0\trans\succ0
  \qquad\text{for every real }\theta\leq\tfrac12.
\end{equation}
All entries of $4N Y_\varepsilon$ are integers in $[0,4N]$.
Every coordinate of $\varepsilon d$ likewise has a positive integer
numerator at most $4N$ when written with denominator $4N$.
Consequently the reduced numerators and denominators of both outputs are
at most $4N$.
\end{lemma}

\begin{proof}
Let $B$ have the generators $u_i,u_g$ as columns, so $M=B\trans B$.
Set
\[
  c^*=(\one_n,\tfrac12\one_p,\xi),\qquad
  \xi=\frac{\eta^2-p/4}{2\eta}=-\frac1{8\eta}.
\]
Their inner products give
\[
  u_i\trans c^*=2+\frac{\abs{S_i}}2=\frac{M_{ii}}2,
  \qquad
  u_g\trans c^*=\frac p4+\eta\xi=\frac{M_{gg}}2.
\]
Thus $\delta=2B\trans c^*$.  The matrix
$\Pi=B(B\trans B)^{-1}B\trans$ is the orthogonal projection onto
$\range(B)$, whence
\[
  s=4\norm{\Pi c^*}^2
    \leq4\norm{c^*}^2
    =4n+p+4\xi^2
    =4n+p+\frac1{4(p-1)}<\frac N2.
\]
For the lower bound, Cauchy--Schwarz in the $M$ inner product gives
\[
  M_{ii}^2
  =(e_i\trans M M^{-1}\delta)^2
  \leq(e_i\trans M e_i)(\delta\trans M^{-1}\delta)
  =M_{ii}s.
\]
Since $M_{ii}>0$, division proves the lower bound.

The binary diagonal equalities follow from $\delta=\diag(M)$.
The lower block $\varepsilon M$ is positive definite, and the Schur
complement of that block in
$Y_\varepsilon-\theta e_0e_0\trans$ is
$1-\theta-\varepsilon s>0$ whenever $\theta\leq1/2$.
This proves \eqref{eq:core-definiteness-margin}.

The entries of $4N Y_\varepsilon$, apart from its root entry $4N$, are
\[
  16+4\abs{S_i},\qquad 2p-1,\qquad
  4\abs{S_i\cap S_j}\ (i\ne j),\qquad 2\abs{S_i}.
\]
They are nonnegative integers.  If a set index exists, $n\geq1$ and
$16+4\abs{S_i}\leq16+4p\leq4N$; all other bounds follow from
$\abs{S_i}\leq p$.  If $n=0$, only the root and $g$ entries remain, and
the bound still holds.  From \eqref{eq:cover-distances}, the numerators of
$\varepsilon d$ with denominator $4N$ are
\[
  16+4\abs{S_i},\qquad 2p-1,\qquad
  32+4\abs{S_i\symdiff S_j}\ (i\ne j),\qquad 2p+15.
\]
These are positive.  Whenever a set index exists,
$4p+32\leq4N$, so they are at most $4N$; without a set index only
$2p-1$ remains.  Reduction of a fraction cannot increase its numerator or
denominator.
\end{proof}

The point $c^*$ is also the sphere center in \eqref{eq:sphere}.
Projecting $c^*$ onto $\range(B)$ leaves the difference of squared
distances in that identity unchanged.
Thus the same sphere that explains the integer slack supplies the bound
on $s$ needed for a scaling with polynomially bounded numerator and
denominator.

\begin{theorem}[Hard points in standard relaxations]
\label{thm:relaxation-hardness}
Suppose the instance is Exact Cover by 3-Sets (X3C):
$\abs{S_i}=3$, $p=3q$, and $q\geq3$.  Let
$\widehat d=\varepsilon d$ and $Y_\varepsilon$ be as above.
\begin{enumerate}
  \item The violated hypermetric inequalities of $\widehat d$ are exactly
  those of
  \[
    b_T=(2-q)e_0+\sum_{i\in T}e_i-e_g,
    \qquad T\text{ an exact cover}.
  \]
  Each has gonality $2q-1$, support size $q+2$, and violation $1/(2N)$.
  All hypermetric inequalities of gonality at most $2q-3$ hold.

  \item The violated rounded positive semidefinite inequalities are exactly
  those described by $b_T$ and $-b_T$; each pair describes the same inequality and has
  violation $1/(2N)$.  No rounded positive semidefinite inequality with
  $\abs{\sigma(b)}\geq3$ is violated.

  \item The violated binary split inequalities at $Y_\varepsilon$ are
  exactly
  \[
    v=(0,-x,1)\quad\text{and}\quad v=(-1,x,-1),
    \qquad x\in\{0,1\}^n,\quad Ax=\one_p.
  \]
  They have $q_{Y_\varepsilon}(v)=-1/(2N)$.

  \item The point $\widehat d$ belongs to $\MET(V)$, and
  $Z(\widehat d)$ belongs to the interior of $\ELL(V)$ relative to its
  affine hull.  Its Gram matrix
  $M(\widehat d)=\varepsilon M$ and binary moment matrix
  $Y_\varepsilon$ are positive definite, and
  $0<\widehat d_{uv}<1$ for distinct $u,v$.
\end{enumerate}
All output numerators and denominators are at most $4N$.
\end{theorem}

\begin{proof}
Summing the equations $Ax=\one_p$ shows that every exact cover uses
exactly $q$ sets.  By \eqref{eq:cover-hypermetric-vectors}, all
hypermetric violators are the stated $b_T$, with gonality
$\abs{2-q}+q+1=2q-1$, support size $q+2$, and
$Q(b_T,\widehat d)=\varepsilon/2=1/(2N)$.

To classify rounded positive semidefinite inequalities, take any integer
vector $b=(b_0,z)$ and write $\sigma=\sigma(b)$.  By \eqref{eq:Q-M}, for
odd $\sigma$ and $u=(\sigma,-2z)$ the rounded positive semidefinite slack at
$\widehat d$ is
\begin{equation}\label{eq:core-rounded-parity}
  \frac{\sigma^2-1}{4}-Q(b,\widehat d)
  =\frac{\sigma^2-1}{4}
    +\varepsilon(z\trans Mz-\sigma\delta\trans z)
  =\frac{u\trans Y_\varepsilon u-1}{4}.
\end{equation}
If $\abs{\sigma}\geq3$, \eqref{eq:core-definiteness-margin} with
$\theta=1/2$ yields
$u\trans Y_\varepsilon u>\sigma^2/2\geq9/2>1$.
If $\sigma=1$, the slack is $\varepsilon g_M(z)$; if $\sigma=-1$,
it is $\varepsilon g_M(-z)$.  Applying \cref{thm:exact-cover} gives
exactly $b_T$ and $-b_T$, with the stated violations.

For a split vector $v$, \eqref{eq:parity-split} expresses its slack in
terms of $2v+e_0$.  This vector has an odd root coordinate and even remaining
coordinates, so \eqref{eq:core-rounded-parity} applies.  The two possible
violated parity vectors are $(1,-2x,2)$ and $(-1,2x,-2)$ for exact covers
$x$, giving precisely the two stated split vectors and their slacks.

It remains to verify the relaxation promises.  Positive definiteness
and the encoding bound follow from \cref{lem:scaling}.  By the congruence
\eqref{eq:congruence}, $Z(\widehat d)$ is positive definite with unit
diagonal.  It lies in the interior of the elliptope.  Its $2\times2$
principal minors are positive, so
$\abs{1-2\widehat d_{uv}}<1$ for distinct $u,v$, proving the strict
coordinate bounds.

Every triangle inequality is a hypermetric inequality of gonality three.
Since $q\geq3$, none is violated.  Every perimeter inequality has vector
$e_u+e_v+e_w$ and sum three, so it holds by the rounded positive
semidefinite classification.
These are the defining inequalities of $\MET(V)$.
\end{proof}

\begin{remark}\label{rem:small-cover-triangles}
The rounded positive semidefinite and split classifications, the positive
definiteness margin, and the elliptope conclusion hold for every exact-cover instance
with $p\geq2$.  In a general instance, all triangle inequalities hold if
and only if no exact cover uses at most two sets.  Indeed, a one-set or
two-set cover in \eqref{eq:cover-hypermetric-vectors} produces a triangle
vector, while every cover with at least three sets produces a vector of
gonality at least five.  For X3C with $q=2$, the hypermetric violators have
gonality three and support size three; these points still satisfy the
elliptope and coordinate promises.
\end{remark}

\begin{corollary}\label{cor:hypermetric-hardness}
Hypermetric violation, rounded positive semidefinite violation, and
Boros--Hammer violation are strongly $\NP$-complete.  For every fixed
$K\geq3$, this remains true
at rational distance vectors $d\in\MET(V)$ with $Z(d)\succ0$ that satisfy
every hypermetric inequality of gonality at most $2K-3$.
The associated Gram and binary moment matrices may both be
required to be positive definite.  The reduction has output numerators
and denominators bounded by $4(8n+2p+1)$, where $n$ is the number of sets
and $p$ the number of elements of the padded X3C instance.
Hypermetric cone membership is strongly $\coNP$-complete.
\end{corollary}

\begin{proof}
X3C is $\NP$-complete \citep{GareyJohnson1979}.  For a fixed $K\geq3$,
add $K$ disjoint triples of new elements and one set for each new triple.
Each added set is forced, so exact covers of the original instance
correspond bijectively to exact covers after padding.  The padded instance
has $q\geq K$, and its size grows by a constant depending only on $K$.
By \cref{thm:relaxation-hardness}, it has an exact cover if and only if
any one of the three inequality families has a violator at the constructed
point.  All the stated relaxation promises hold on both yes- and
no-instances.

The source instance has an explicit finite universe.  Computing the
intersections in \eqref{eq:cover-matrix} takes polynomial time, and the
output has $O((n+2)^2)$ entries with numerators and denominators
$O(n+p)$.  Thus its encoding is polynomial even if the numerical data
are encoded in unary, proving strong $\NP$-hardness.

The short-certificate results
\cref{lem:short-hypermetric-witness,lem:short-split-witness} prove
$\NP$ membership for arbitrary rational inputs.  For any rational distance
vector, set $M=M(d)$, $\delta=\diag(M)$, and $Y=Y(\delta,M)$.
The algebra in \eqref{eq:core-rounded-parity} applies with
$\varepsilon=1$ and identifies a rounded positive semidefinite vector
of odd sum with a split vector: $u=(\sigma,-2z)=2v+e_0$ and
$v=((\sigma-1)/2,-z)$.  This transformation preserves polynomial
certificate length in both directions.  Boros--Hammer inequalities are
the binary moment form of those split inequalities.  Taking the complement
of hypermetric violation proves the cone-membership statement.

For hypermetric violation alone, homogeneity also permits the integral
output $4d$ in \eqref{eq:cover-distances}.  Its coordinates are at most
$4p+32$, and its Gram matrix is positive definite.  It satisfies all
triangle inequalities and the same gonality promise, so this also gives
strong hardness with integral distances.
\end{proof}
